% XIII LAW3M 2026 abstract
% Prepared according to the abstract template: one page, Times, 12 pt body, up to three references.
\documentclass[12pt,a4paper]{article}

\usepackage[utf8]{inputenc}
\usepackage[T1]{fontenc}
\usepackage[english]{babel}
\usepackage{mathptmx}
\usepackage{geometry}
\usepackage{setspace}
\usepackage[normalem]{ulem}
\usepackage{microtype}
\geometry{top=2.0cm,bottom=2.0cm,left=2.2cm,right=2.2cm}
\pagestyle{empty}
\setlength{\parindent}{0pt}
\setlength{\parskip}{0pt}
\singlespacing

\begin{document}

\begin{center}
{\bfseries\fontsize{14}{16}\selectfont XIII LAW3M - Abstract Submission Template}\par
\vspace{6pt}
{\bfseries\fontsize{14}{16}\selectfont dipolar compass-needle arrays as a minimal model for static and dynamic collective magnetism}\par
\vspace{10pt}
{\bfseries\fontsize{12}{14}\selectfont \uline{Jo\~ao~P.~Sinnecker}}\par
\vspace{6pt}
{\itshape\fontsize{12}{14}\selectfont Centro Brasileiro de Pesquisas F\'isicas (CBPF), Rio de Janeiro, Brazil}\par
\end{center}
\vspace{10pt}

Artificial arrays of interacting nanomagnets provide a direct route to study how local dipolar coupling produces collective magnetic states. However, in nanolithographic systems it is often difficult to separate the effects of geometry, spacing, disorder, damping, and field history. Here we use a two-dimensional array of magnetic compass needles as a minimal physical model for this problem. Each needle is treated as a point magnetic dipole fixed at a lattice site and free to rotate in the plane under the total dipolar field generated by all other needles, an external magnetic field, and viscous damping. The model keeps the full long-range and anisotropic dipolar interaction and follows the Newtonian rotational dynamics of each element. This approach is intentionally simple, but it captures a central question in artificial magnetic materials: how do collective configurations emerge when many nearly bistable magnetic units interact mainly through their stray fields? By varying the lattice geometry, needle spacing, and damping, the simulations allow us to explore regimes that are difficult to isolate experimentally in nanofabricated arrays. Square, triangular, and honeycomb arrangements are compared as controlled examples of different interaction topologies. The simulations reveal how dipolar coupling can promote domain formation, cooperative reversal, hysteresis, and avalanche-like changes in the magnetic configuration. Preliminary results indicate that the avalanche statistics and the hysteresis response are not fixed properties of the lattice alone, but depend strongly on the dynamical regime set by the damping. The compass-needle array should therefore not be viewed as a literal replacement for nanomagnetic islands. Its value is as a transparent toy model in which the role of dipolar interactions can be examined without the additional complications of fabrication defects, edge roughness, and material-dependent switching barriers. In this sense, it provides a useful bridge between classroom-scale magnetic analogues and current problems in artificial spin systems, especially the relation between lattice design, dipolar frustration, and magnetic dynamics.

{\small
[1] A.~Clauset, C.~R.~Shalizi and M.~E.~J.~Newman, SIAM Rev.\ \textbf{51}, 661 (2009).\par
[2] R.~Taniguchi, Sci.\ Rep.\ \textbf{15}, 12345 (2025).\par
[3] K.~A.~Dahmen, Y.~Ben-Zion and J.~T.~Uhl, Nat.\ Phys.\ \textbf{7}, 554 (2011).}

\end{document}

